% Generated draft: requires mathematical and source review.
\documentclass[11pt]{article}
\usepackage[margin=25mm]{geometry}
\usepackage{fontspec}
\setmainfont{Latin Modern Roman}
\usepackage{amsmath,amssymb,mathtools,unicode-math}
\setmathfont{Latin Modern Math}
\usepackage{enumitem,needspace,xurl,hyperref}
\setlength{\parindent}{0pt}
\setlength{\parskip}{5pt}
\setlength{\emergencystretch}{3em}
\newcommand{\N}{\mathbb{N}}
\newcommand{\Z}{\mathbb{Z}}
\newcommand{\Q}{\mathbb{Q}}
\newcommand{\R}{\mathbb{R}}
\newcommand{\C}{\mathbb{C}}
\newcommand{\F}{\mathbb{F}}
\newcommand{\A}{\mathbb{A}}
\newcommand{\PP}{\mathbb{P}}
\newcommand{\E}{\mathbb{E}}
\newcommand{\eps}{\varepsilon}
\newcommand{\dd}{\,\mathrm{d}}
\begin{document}
\section*{UnsolvedMath ID: 10400074}
\textbf{Source problem code:} catalogue entry
\subsection{Definitions and mathematical statement}
Let \(\mathcal A_\Sigma\) be the Poisson algebra and \(\mathcal M_G(\Sigma)\) the moduli space referred to in the source. The surface \(\Sigma\), group \(G\), algebra generators and relations, and function class in \(C(\mathcal M_G(\Sigma))\) are not defined in this extract. Retain those symbols with their unresolved source meanings.

A formal deformation quantization of a complex Poisson algebra \((A,\cdot,\{\ ,\ \})\), in one common normalization, is an associative \(\C[[\hbar]]\)-bilinear product on \(A[[\hbar]]\) of the form
\[
 a*b=ab+\sum_{r\ge1}\hbar^r B_r(a,b),\qquad
 a*b-b*a=\hbar\{a,b\}+O(\hbar^2),
\]
with appropriate unit and continuity conditions. A convention using \(i\hbar\) instead must be fixed by the source, not inferred here.

T. Kohno's Problem 3.24 asks for a deformation quantization of \(\mathcal A_\Sigma\) which descends to a deformation quantization of \(C(\mathcal M_G(\Sigma))\). If descent is through a source-defined Poisson quotient map \(q:\mathcal A_\Sigma\to C(\mathcal M_G(\Sigma))\), a precise formal version asks for products \(*\) and \(*_{\mathcal M}\) and a deformation \(q_\hbar\) of \(q\) satisfying
\[
 q_\hbar(a*b)=q_\hbar(a)*_{\mathcal M}q_\hbar(b),
 \qquad q_\hbar\bmod\hbar=q.
\]
In a quotient realization, the deformed relations must define an appropriate two-sided ideal and the required formal quotient, rather than merely a product on the larger algebra. The record supplies neither \(q\) nor the quotient relations, so this is a conditional explanation of descent, not a specified new quotient. The source additionally asks to investigate the relation with geometric quantization of \(\mathcal M_G(\Sigma)\); it gives no proposed comparison map or equivalence criterion.
\subsection{Short English statement}
Quantize the source's Poisson algebra for a surface in a way that passes to its stated moduli-space function algebra, and investigate the relation with geometric quantization. The moduli problem, function regularity, and descent map are missing. The supplied partial-progress summary mentions extensive quantization literature without identifying a result that settles this exact request.
\subsection{Sources}
\begin{itemize}
\item Local UnsolvedMath registry, permanent ID \(10400074\), exact source code \texttt{catalogue entry}, Problem 3.24, attributed to T. Kohno. The supplied title is truncated after ``(T.''
\item Ohtsuki, \emph{Problems on invariants of knots and 3-manifolds} (2002), Problem 3.24, PDF page 71: \url{https://msp.org/gtm/2002/04/gtm-2002-04-024s.pdf}. Printed bold headings distinguish 231 numbered items from references in exposition.
\item Extraction pdf-font-structure, accessed 2026-07-29. Source evidence NEEDS\_REVIEW presents an open problem and requires a current primary-literature audit. Registry status partially\_solved and research classification PARTIAL-PROGRESS mention quantization of moduli spaces and character varieties, with no further specific references or precise resolved subcase.
\item Dataset attribution: Ulam AI / UnsolvedMath Contributors, CC BY 4.0. Source rights NEEDS\_REVIEW: publicly accessible author/publisher copy, redistribution terms requiring review.
\end{itemize}
\end{document}
